\documentclass[12pt]{report}

%\usepackage[utf8]{inputenc} 			% utf-8, xetex es utf-8 por defecto
\usepackage{etoolbox}					% quitar espacio a lista de figuras %
\usepackage{tabularx,booktabs}  			% uso de tablas %
\usepackage[labelfont=bf]{caption}			% para figuras y tablas usar la bajada
\usepackage[spanish]{babel}				% idioma español para  %}
\usepackage{comment} 					% para el \begin{comment}
\usepackage{rotating} 					% rotar las tablas grandes de los apéndices %
\usepackage{longtable}					% tablas de varias páginas, hecha para los apendices %
\usepackage{lscape}					% tablas de varias páginas, hecha para los apendices %
\usepackage{graphicx}					% para incluir imagenes %
\usepackage{apacite}					% citas de tipo APA 6th %
\usepackage[onehalfspacing]{setspace}		% interlineado 1.5 para el texto%
\usepackage{tocloft}					% para darle estilo a la tabla de contenidos %
\usepackage{multirow}					% para unir multiples filas en una tabla %
\usepackage[titletoc]{appendix}			% agregar anexos
\usepackage{fancyhdr}					% encabezados y pies de pagina
\usepackage[explicit]{titlesec}				% quitar  titulos de secciones
\usepackage{fontspec}					% cambiar fuente
\usepackage[includeheadfoot]{geometry}	% cambiar dimensiones pagina
\usepackage{colortbl}					% colorear tabla
\usepackage{url}

%% quitar guiones cuando las palabras sobrepasan el reglón
\tolerance=1
\emergencystretch=\maxdimen
\hyphenpenalty=10000
\hbadness=10000



% margenes y tamaños
\geometry{
 letterpaper,
 total={216mm,279mm},
 left=30mm,
 right=20mm,
 top=12mm,
 bottom=14mm,
}


% tabs para la seccion de comision examinadora
\newcommand{\tab}[1]{\hspace{.3\textwidth}\rlap{#1}}
\newcommand{\itab}[1]{\hspace{0em}\rlap{#1}}

%% esconder los titulos
\newcommand*\Hide{%
\titleformat{\chapter}[display]
  {}{}{0pt}{\Huge}
\titleformat{\part}
  {}{}{0pt}{}
}

% titulares de cada capitulo en formato
\newenvironment{titular}{%
	\thispagestyle{empty}
	\setlength{\parindent}{0cm}
	\fontsize{22pt}{26.4pt}%
	\selectfont%
    	\center%
    	\begingroup%
}{%
   	\endgroup
    	\endcenter
	\clearpage
	\par
}


%%%%%%%%% Vars %%%%%%%%%
\newcommand{\titulotesis}{Titulo de la tesis}
\newcommand{\nombre}{JAVIER IGNACIO FUENTES MUÑOZ}
\newcommand{\anyo}{2016}


%%%%%%%% END VARS %%%%%%

% encabezados y pie de pagina
\renewcommand{\headrulewidth}{0pt}
\renewcommand{\footrulewidth}{0pt}
\setlength{\footskip}{25pt}
\lhead{}
\chead{}
\renewcommand{\chaptermark}[1]{\markboth{\thechapter.\space#1}{}} 
\lfoot{\fontsize{10pt}{12pt}\selectfont \textbf{\titulotesis}}

\cfoot{}

\rfoot{\thepage}

\fancypagestyle{plain}{\fancyhf{}
	\fancyfoot[LE,LO]{\fontsize{10pt}{12pt}\selectfont \textbf{\titulotesis}}
	\fancyfoot[RE,RO]{\thepage}
	\fancyhead[RE,RO]{\it\rightmark}
	\renewcommand{\headrulewidth}{0ex}
} 

%\counterwithout{table}{chapter}	% para contar los cuadros en 1,2,3 . . . en vez de 1.1,1.2,1.3,2.1, ... %
%\counterwithout{figure}{chapter}	% para contar las figuras en 1,2,3 . . . en vez de 1.1,1.2,1.3,2.1, ... %
%\tolerance=10000			% evita guiones que separen las palabras %
\graphicspath{{./img/}} 			% directorio de imagenes &
\renewcommand\spanishtablename{Tabla\ } 		% cambiar caption de los cuadrs


% modificar los indices
\AtBeginDocument{\renewcommand\appendixname{Anexo}}			%cambiar nombre en el indice
\AtBeginDocument{\renewcommand\listtablename{Índice de tablas}}		%cambiar nombre en el indice
\renewcommand{\cftfigpresnum}{Figura\ }
\renewcommand{\cfttabpresnum}{Tabla\ }
% espaciado en el prefijo del indice de las tablas "Figura X .." y "Tabla X ..."
\newlength{\mylenf}
\settowidth{\mylenf}{\cftfigpresnum}
\setlength{\cftfignumwidth}{\dimexpr\mylenf+3.5em}
\setlength{\cfttabnumwidth}{\dimexpr\mylenf+3.5em}

% quitar espacio anterior tabla de contenidos,figuras y tablas %
\setlength{\cftbeforetoctitleskip}{-3em}
\setlength{\cftbeforeloftitleskip}{-3em}
\setlength{\cftbeforelottitleskip}{-3em}

% quitar espacio entre capitulos en la lista de figuras y tablas%
\makeatletter
\patchcmd{\@chapter}{\addtocontents{lof}{\protect\addvspace{10\p@}}}{}{}{}% LoF
\patchcmd{\@chapter}{\addtocontents{lot}{\protect\addvspace{10\p@}}}{}{}{}% LoT
\makeatother


\begin{document}

% la cita de multiples autores es con "et al." en vez de "y cols"
\renewcommand{\BOthers}[1]{et al.\hbox{}}


\singlespace				% interlineado 1.0
\pagenumbering{gobble}		% sin numerar las primeras paginas %
\pagestyle{empty}			% sin encabezados

\leftskip=0cm 
\rightskip=0cm

\begin{center}
\begin{tabularx}{\textwidth}{p{3cm} >{\centering\arraybackslash}X}
	\vspace{0pt} 
	\includegraphics[width=2.8cm]{logo}
   	& 
   	\vspace{10pt} \textbf{ \uppercase{universidad de la frontera \linebreak facultad de ingeniería y ciencias\linebreak depto. cs. computación e informática}}
	
\end{tabularx}
\end{center}

\null
\vfill

\begin{center}
	\textbf{ \uppercase{
	``\titulotesis''
	} }
\end{center}

\null
\vfill

\begin{center}
	\textbf{ \uppercase{
	\nombre \\
	\anyo
	} }
\end{center}

\clearpage
\newpage\null\thispagestyle{empty}\newpage

\begin{center}
\begin{tabularx}{\textwidth}{p{3cm} >{\centering\arraybackslash}X}
	\vspace{0pt} 
	\includegraphics[width=2.8cm]{logo}
   	& 
   	\vspace{10pt} \textbf{ \uppercase{universidad de la frontera \linebreak facultad de ingeniería y ciencias\linebreak depto. cs. computación e informática}}
	
\end{tabularx}
\end{center}

\null
\vfill

\begin{center}
	\textbf{ \uppercase{
	``\titulotesis''
	} }
\end{center}

\null
\vfill

\rightskip=0cm
\noindent\rule{\textwidth}{1.5pt}
\setlength{\parindent}{0cm}
{
	\raggedright

	\uppercase{\textbf{trabajo para optar al título \\de ingeniero civil industrial mención informática}
}
\noindent\rule{\textwidth}{1.5pt}

\vspace{15 mm}
\leftskip=0cm

\uppercase{\textbf{\hfill Profesor guía:} \profesorguia}\\
\vspace*{0.5cm}
\uppercase{\textbf{\hfill Profesor co-guía:} \profesorcoguia}
}


\null
\vfill

\begin{center}
	\textbf{ \uppercase{\nombre \\ \anyo}}
\end{center}

\clearpage

\leftskip=0cm
\rightskip=0cm
\setlength{\parindent}{0cm}

\begin{center}
	\textbf{ \uppercase{
	``\titulotesis'' \\
	\vspace{5 mm}
	\nombre \\
	\vspace{5 mm}
	comisión examinadora
	} }
	\vspace{25 mm}

	\textbf{ \uppercase{\profesorguia} \\
		Profesor Guía
	}
	
	\vspace*{1.5cm}
		
	\textbf{ \uppercase{\profesorcoguia} \\
		Profesor Co-Guía
	}

\end{center}

\vspace{4cm}

\begin{tabularx}{\textwidth}{ X X X } 
	Dr. CARLOS FERNANDO CARES GALLARDO & Mg. JORGE ALBERTO HOCHSTETTER DIEZ  & Otro profesor connotado\\
	\multicolumn{1}{c}{Profesor Examinador 1} &	\multicolumn{1}{c}{Profesor Examinador 2} & \multicolumn{1}{c}{Profesor Examinador 3} \\					
\end {tabularx}
\vfill
\begin{flushright}
	\textbf{Nota trabajo escrito :} \tab{} \\ 	
	\textbf{Nota examen :} \tab{} \\ 
	\textbf{Nota final :} \tab{} \\ 		
\end{flushright}

			

\onehalfspacing			% interlineado 1.5


{
\fontsize{16pt}{19.2pt}%
\selectfont%
\center%
\begingroup%

%\uppercase{
%dedicatoria 
%}
\endgroup
\endcenter
}

{
\topskip0pt
\vspace*{\fill}
\hfill \textit{A mi abuela Joaquina}
\vspace*{\fill}
}

\clearpage

\begin{comment}
{
\fontsize{16pt}{19.2pt}%
\selectfont%
\center%
\begingroup%

\uppercase{
	agradecimientos 
}
	

\endgroup
\endcenter
}


\clearpage

\end{comment}

{
\fontsize{16pt}{19.2pt}%
\selectfont%
\center%
\begingroup%

\uppercase{
	resumen 
}	

\endgroup
\endcenter
}

\setlength{\parindent}{1cm}
\setlength{\parskip}{5pt}

%%%%%
Resumen here

\clearpage                       % Otherwise \pagestyle affects the previous page.
{                                % Enclosed in braces so that re-definition is temporary.
  \pagestyle{empty}              % Removes numbers from middle pages.
  \fancypagestyle{plain}         % Re-definition removes numbers from first page.
  {
    \fancyhf{}%                       % Clear all header and footer fields.
    \renewcommand{\headrulewidth}{0pt}% Clear rules (remove these two lines if not desired).
    \renewcommand{\footrulewidth}{0pt}%
  }
  \tableofcontents
  \newpage
  \listoffigures
  \newpage
  \listoftables
  \thispagestyle{empty}          % Removes numbers from last page.
}

% estilos para el documento %
\setlength{\parindent}{1cm}
\setlength{\parskip}{5pt}
\renewcommand{\arraystretch}{1.5}

\pagestyle{fancy}
\pagenumbering{arabic} 		% empezar la numeracion	

{
\Hide
\chapter{Introducción}
}

\begin{titular} 
	\uppercase{
	capítulo 1 \\
	introducción \\
	}
\end{titular}

\section{Descripción del problema}

\subsection{Exposición general del problema}
ACS es un software distribuido altamente complejo. Como toda infraestructura de software, este presenta fallas en su funcionamiento las que pueden llegar a provocar la detención completa del sistema, provocando retrasos en las tareas de observación.
Gran parte del funcionamiento de ACS es permanentemente registrado en archivos de registros, los cuales suman diariamente grandes cantidades de información sobre la operación y eventos asociados a diferentes niveles de alerta. Sin embargo, dada la gran cantidad de registros generados por este sistema, resulta complejo el poder analizar estos datos en forma eficiente y efectiva.
 
ALMA tiene una infraestructura altamente compleja, lo que se refleja en sus sistema de generación de logs, por lo que es natural encontrar problemáticas asociadas a numerosas causas y por este motivo es necesario buscar y analizar estas posibles problemáticas. El problema se plantea complejo tanto por la diversidad de los logs como por su cantidad.
 
Por otra parte, se ha desarrollado herramientas y metodologías de minería de datos que tienen por objetivo el apoyar el análisis de grandes volúmenes de datos de manera de poder extraer conocimientos nuevos y útiles. Estas herramientas son capaces de procesar grandes volúmenes de información y encontrar relaciones entre las múltiples variables que afectan los procesos. Para poder aplicar estas herramientas es necesario plantear objetivos de búsqueda y aplicar una metodología de minería de datos CRSIP.
 
Sin embargo, debido a la alta complejidad del sistema ACS, los objetivos de búsqueda que guían el análisis son múltiples, lo que requiere de una comprensión muy acabada de todo el proceso para poder modelar algunas partes a través de la minería de datos y buscar entonces las relaciones que permiten comprender por ejemplo, porqué y bajo qué circunstancias el sistema falla. El estudio y comprensión del sistema ACS permitirá la identificación de un problema específico para el cual es posible plantear una aproximación basada en minería de datos.

\subsection{Situación actual del problema}
Actualmente los registros de ACS son numerosos, el origne de estos registros provienen en su mayoria de todos los sistemas y subsistemas, desde operaciones de alto nivel hasta operaciones de bajo nivel. Todos estos registros son almacenados en un mismo formato y lugar. Para acceder a estos registros es posible a traves del servicio de Kibana, que permite realizar búsquedas y obtener estos registros.
La idea principal es reducir el espacio generado por los registros de ACS en el contexto de un proceso de observación.
Para este enfoque se plantea por parte de ALMA utilizar el concepto de máquinas de estado finitas (MEF) para reducir el conjunto de datos a estudiar, generando un modelo sobre el proceso de observación. 
\section{Objetivos}
\subsection{General}
\begin{itemize}
	\item Desarrollar una herramienta de pre-procesamiento para mineria de datos basada en los registros de ACS.
\end{itemize}
\subsection{Específicos}
\begin{itemize}
	\item 
\end{itemize}

{

\Hide
\chapter{Antecedentes Generales}
}

\begin{titular} 
	\uppercase{
	capítulo 2 \\
	Antecedentes Generales \\
	}
\end{titular}


{
\Hide
\chapter{Descripción de Actividades Realizadas}
}

\begin{titular} 
	\uppercase{
	capítulo 3 \\
	Descripción de Actividades Realizadas \\
	}
\end{titular}





{

\Hide
\chapter{Resultados y discusión}
}

\begin{titular} 
	\uppercase{
	capítulo 4 \\
	Resultados y discusión \\
	}
\end{titular}



{
\Hide
\chapter{Conclusiones}
}

\begin{titular} 
	\uppercase{
	capítulo 5 \\
	Conclusiones \\
	}
\end{titular}

{
	\fontsize{22pt}{26.4pt}%
	\selectfont%
	Conclusiones
}

 




%%% corregir posicion del titulo de las referencias
\makeatletter
\def\@makechapterhead#1{%
  \vspace*{-40\p@}%
  {\parindent \z@ \raggedright \normalfont
    \ifnum \c@secnumdepth >\m@ne
        \huge\bfseries \@chapapp\space \thechapter
        \par\nobreak
        \vskip 20\p@
    \fi
    \interlinepenalty\@M
    \Huge \bfseries #1\par\nobreak
    \vskip 2\p@
  }}
\def\@makeschapterhead#1{%
  \vspace*{-40\p@}%
  {\parindent \z@ \raggedright
    \normalfont
    \interlinepenalty\@M
    \Huge \bfseries  #1\par\nobreak
    \vskip 2\p@
  }}
\makeatother


\pagestyle{fancy}
\renewcommand{\spanishbibname}{Bibliografía} 	% cambiar de nombre a la referencia
\bibliographystyle{apacite}
\bibliography{Tesis}

\appendix
\begin{appendices}

{
\Hide
\chapter{Tabla de Puntajes}
}

\begin{titular} 
	\uppercase{
	Anexo A \\
	Tabla de Puntajes \\
	}
\end{titular}




{
\Hide
\chapter{Lista de publicaciones}
}

\begin{titular} 
	\uppercase{
	Anexo B \\
	Lista de publicaciones \\
	}
\end{titular}


%\include{apendice2-codigo}
\end{appendices}


\end{document}

